% DRAFT subsection for the Results -- "attack-unlimited" / voracious mixture.
% Numbers are from the converged Model 2a posterior (outputs/model_2a_posterior.csv,
% outputs/model_2a_bayes_pvalues.csv, outputs/loo_comparison.csv). Cross-references
% (Fig./Table/Supp.) are placeholders to be wired to the manuscript.

\subsection{Prey searching is negligible for most fish}

Posterior predictive checks revealed one systematic failure of the two-state
model: it substantially under-predicted the number of fish that consumed
\emph{all} of their available prey. Across prey densities, 29\% of fish cleared
their tanks, whereas the two-state model predicted only $\approx$8\%
(Supplementary Fig.~S\ref{supp:ppc}). This ``full-consumption'' discrepancy was
strongest among heavier individuals.

To accommodate this behaviour, we extended the model to include a latent
subpopulation of \emph{voracious} foragers, for whom prey capture is effectively
instantaneous---formally, the limit $a \to \infty$ of eq.~(\ref{eqn:eqn1}). In
this limit a fish is constrained only by prey handling, and the number of prey
eaten during a trial follows a shifted, right-truncated Poisson distribution,
$\min\!\big(N_0,\, 1 + \mathrm{Poisson}(r)\big)$ (derivation in
\hyperlink{supp:S1}{Supplementary File S1}). Each fish belongs to the voracious
class with a weight-dependent probability $\pi$, where
$\mathrm{logit}(\pi) = \gamma_0 + \gamma_w\,\mathrm{weight}$, and the latent class
membership is marginalised out of the likelihood. This mixture resolved the
full-consumption misfit (posterior-predictive $p = 0.73$;
Supplementary Fig.~S\ref{supp:ppc}) and was decisively favoured over the
two-state model by leave-one-out cross-validation, which also became reliable
only under the mixture (Table~\ref{tab:model_comparison}).

Strikingly, the fitted model estimated that the great majority of fish---92\%
(95\% credible interval: 90--95\%)---belonged to the voracious class, and this
proportion increased with body weight, from $\approx$78\% for the lightest fish
to $\approx$100\% for the heaviest (Fig.~\ref{fig:voracious}). In other words,
for almost all individuals the time spent searching for prey was negligible
relative to handling time, so that consumption was governed almost entirely by
the rate of prey handling, $r$. This provides a quantitative basis for the
observation that searching time varied little across treatments: within the
confined experimental arenas, prey encounter was so rapid that the searching
phase contributed negligibly to feeding for the overwhelming majority of fish.
Importantly, allowing for this subpopulation did not change our conclusions about
how the treatments affected feeding, because the treatment covariates act on the
handling rate $r$, which is shared by both components of the mixture.
